\begin{notechapter}
\chapter{Integration on Manifolds and Stokes' Theorem}
\label{ch:stokes}
\chaptersubtitle{%
  The fundamental theorem of calculus in every dimension}

\begin{chapterintro}
The integral of \cref{ch:lebesgue-integral} lives on $\R^n$. To
integrate on a manifold one works in charts, and the result must not
depend on the charts. The change of variables theorem says which
objects have this property: those that pick up the factor
$\abs{\det}$ of the Jacobian matrix, called densities. On an oriented
manifold the forms of top degree are densities, and for them there is
one more theorem, due to Stokes in its general form: the integral of
$\dd\omega$ over a domain is the integral of $\omega$ over its
boundary. It contains the fundamental theorem of calculus and all the
integral theorems of vector calculus. You will learn
\begin{itemize}[nosep]
  \item what an orientation of a manifold is, and that it may not exist;
  \item densities, their integral, and the volume of a Riemannian
        manifold;
  \item the integral of a $k$-form over an oriented manifold of
        dimension $k$;
  \item domains with boundary, the orientation of the boundary, and how
        to cut a form into pieces that live in one chart;
  \item Stokes' theorem, with its proof and first consequences;
  \item the same theorem on a graph, where the boundary is the dual of
        the differential;
  \item integration by parts as the source of the identities behind
        score matching.
\end{itemize}
\prerequisites{\cref{ch:tangent-bundles} for manifolds;
\cref{ch:lebesgue-integral,ch:fubini,ch:change-of-variables} for the
integral, Fubini's theorem and the change of variables;
\cref{ch:multilinear-algebra,ch:differential-forms} for orientation,
forms, pull-back and $\dd$. Throughout, $M$ is a manifold of dimension
$n$. Its charts are written $(U, \varphi)$ or $(U, x)$, with
$x = (x_1,\dots,x_n)$ the coordinate functions.}
\end{chapterintro}

\section{Orientation}\label{stokes:sec:orientation}

\begin{definition}[Orientation of a manifold]\label{stokes:def:orientation}
An \emph{orientation} of $M$ is the choice of an orientation of each
tangent space $T_pM$ (\cref{multilinear-algebra:def:orientation}) such
that $M$ is covered by \emph{positive charts}: charts $(U, x)$ whose
coordinate basis $(\partial_1,\dots,\partial_n)$ is positive at every
point of $U$. A manifold is \emph{orientable} if it has an orientation,
and \emph{oriented} when one is chosen.
\end{definition}

On the overlap of two charts the coordinate bases are related by the
Jacobian matrix $J$ of the transition map
(\cref{tangent-bundles:thm:change}). So two positive charts have
$\det J > 0$ on their overlap. Conversely, an atlas all of whose
transition maps have $\det J > 0$ defines an orientation: declare the
coordinate bases of its charts positive.

\begin{proposition}[Orientation by a form]\label{stokes:prop:form}
If an $n$-form $\mu$ on $M$ vanishes at no point, the bases on which
$\mu$ is positive define an orientation of $M$.
\end{proposition}

\begin{proof}
Let $p\in M$ and $(U, x)$ a chart around $p$, with
$\mu = f\,\dd x_1\wedge\dots\wedge\dd x_n$. The function $f$ is
continuous and $f(p)\neq0$, so $f$ keeps its sign on a smaller open set
around $p$, to which we restrict the chart. If $f < 0$ there, replace
$x_1$ by $-x_1$. The chart is now positive.
\end{proof}

\begin{example}\label{stokes:ex:orientations}
\begin{enumerate}[nosep]
  \item An open subset of an oriented affine space is oriented by the
        affine charts of the positive frames; $\R^n$ is oriented by its
        canonical basis.
  \item \emph{Hypersurfaces.} Let $\Aff$ be an oriented Euclidean
        affine space and $M = \{g = 0\}$, where $\dd g\neq0$ at the
        points of $M$. The gradient $\nabla g$ is not tangent to $M$,
        since $\dd g(\nabla g) = \norm{\nabla g}^2 > 0$. So for a basis
        $(v_2,\dots,v_n)$ of $T_pM$ the number
        $\operatorname{vol}(\nabla g, v_2,\dots,v_n)$ is not $0$: the
        restriction to $M$ of $i_{\nabla g}\operatorname{vol}$ vanishes
        nowhere and orients $M$. A basis of $T_pM$ is positive when it
        is positive after the normal vector $\nabla g(p)$ is put in
        front. The sphere $\norm x^2 = 1$ is oriented in this way by
        its outward normal.
  \item \emph{The M\"obius band} is the quotient of the strip
        $S = \R\times(-1, 1)$ by $(s, t)\sim(s + 1, -t)$, with the
        charts given by the pieces $(a, a + 1)\times(-1, 1)$ of the
        strip. It is not orientable. An orientation would lift to an
        orientation of $S$ whose sign with respect to the canonical one
        is locally constant, hence constant along every segment, hence
        constant on $S$. It would also be preserved by the gluing map
        $(s, t)\mapsto(s + 1, -t)$, whose determinant is $-1$: a
        contradiction (\cref{stokes:fig:moebius}).
\end{enumerate}
\end{example}

\begin{figure}[ht]
\centering
\begin{tikzpicture}[>=Stealth, scale=1.0]
  \foreach \sh/\dir/\lab in {0/1/cylinder, 7.4/-1/M\"obius band} {
    \begin{scope}[shift={(\sh,0)}]
      \fill[bookintuition!14] (0,0) rectangle (5,1.8);
      \draw[bookintuition!70!black, dashed] (0,0) -- (5,0);
      \draw[bookintuition!70!black, dashed] (0,1.8) -- (5,1.8);
      \draw[very thick, bookaccent, ->] (0,0.2) -- (0,1.6);
      \ifnum\dir=1
        \draw[very thick, bookaccent, ->] (5,0.2) -- (5,1.6);
      \else
        \draw[very thick, bookaccent, ->] (5,1.6) -- (5,0.2);
      \fi
      \foreach \x in {0.9,2.5} {
        \draw[->, thick] (\x,0.9) -- ++(0.55,0);
        \draw[->, thick, bookpitfall] (\x,0.9) -- ++(0,0.55);
      }
      \draw[->, thick] (4.0,0.9) -- ++(0.55,0);
      \draw[->, thick, bookpitfall] (4.0,0.9) -- ++(0,0.55);
      \draw[->, thick] (0.25,0.9) -- ++(0.4,0);
      \draw[->, thick, bookpitfall] (0.25,0.9) -- ++(0,{\dir*0.45});
      \node at (2.5,-0.45) {\lab};
    \end{scope}
  }
\end{tikzpicture}
\caption{A strip whose two ends are glued as the arrows indicate. A
basis carried once around the cylinder comes back as it was. On the
M\"obius band it comes back with its second vector reversed, at the far
left of the picture: no choice of positive bases is continuous.}
\label{stokes:fig:moebius}
\end{figure}

\section{Densities}\label{stokes:sec:densities}

A function on $M$ cannot be integrated by integrating its expression in
a chart: the result changes with the chart, by the change of variables
theorem. The objects that can be integrated carry the compensating
factor.

\begin{definition}[Density]\label{stokes:def:density}
A \emph{density} $\rho$ on $M$ assigns to each chart $(U, \varphi)$ a
Borel function $\rho_\varphi:\varphi(U)\to\R$, such that for two charts
with transition map $\tau = \varphi'\circ\varphi^{-1}$,
\begin{equation}\label{stokes:eq:density}
  \rho_\varphi = (\rho_{\varphi'}\circ\tau)\,\abs{\det D\tau}
  \qquad\text{on }\varphi(U\cap U') .
\end{equation}
It is continuous, smooth or non-negative if the functions
$\rho_\varphi$ are.
\end{definition}

It is enough to give the functions $\rho_\varphi$ for the charts of one
atlas, with \eqref{stokes:eq:density} on the overlaps: the rule then
defines $\rho_\varphi$ for every other chart, consistently, because
determinants multiply under composition. Densities are added and
multiplied by functions, $(h\rho)_\varphi
= (h\circ\varphi^{-1})\rho_\varphi$, and $\abs\rho$ is a density.

\begin{example}\label{stokes:ex:densities}
\begin{enumerate}[nosep]
  \item \emph{Riemannian volume.} On a Riemannian manifold with Gram
        matrix $G$ in the chart $\varphi$
        (\cref{riemannian-optimization:def:metric}), put
        $\rho_\varphi = \sqrt{\det G}\circ\varphi^{-1}$. In another
        chart $G = J^\top G'J$, so
        $\sqrt{\det G} = \sqrt{\det G'}\,\abs{\det J}$: this is a
        density, the \emph{volume density}. For a submanifold of a
        Euclidean space with the induced metric it is the density of
        \cref{change-of-variables:def:volume}.
  \item \emph{Forms.} For an $n$-form
        $\omega = f\,\dd x_1\wedge\dots\wedge\dd x_n$, the functions
        $\abs f\circ\varphi^{-1}$ are a density $\abs\omega$, by
        \cref{differential-forms:thm:top}. If $M$ is oriented, the
        functions $f\circ\varphi^{-1}$ themselves, for the positive
        charts, are a density $\rho^\omega$: between positive charts
        $\det D\tau = \abs{\det D\tau}$.
\end{enumerate}
\end{example}

A subset $B$ of $M$ is \emph{Borel} if $\varphi(B\cap U)$ is a Borel
subset of $\R^n$ for every chart, and \emph{null} if these sets are
Lebesgue-null. Both conditions need only be checked on a countable
atlas, because transition maps are diffeomorphisms
(\cref{change-of-variables:cor:sets}). Let now $(U_i, \varphi_i)$,
$i = 1, 2,\dots$, be a countable atlas
(\cref{tangent-bundles:def:manifold}) and cut $M$ into the disjoint
Borel pieces
\[
  E_1 = U_1 , \qquad E_i = U_i\setminus(U_1\cup\dots\cup U_{i-1}) .
\]

\begin{theorem}[Integral of a density]\label{stokes:thm:density}
Let $\rho\geq0$ be a density and $B\subset M$ a Borel set. The number
\[
  \int_B\rho = \sum_i\int_{\varphi_i(B\cap E_i)}\rho_{\varphi_i}\,
  \dd\lambda\ \in[0, \infty]
\]
does not depend on the countable atlas or on its ordering. It equals
$\int_{\varphi(B)}\rho_\varphi\,\dd\lambda$ when $B$ lies in the domain
of a chart $\varphi$. The map $B\mapsto\int_B\rho$ is a measure on the
Borel sets of $M$, for which null sets have measure $0$.
\end{theorem}

\begin{proof}
Let $(U'_k, \varphi'_k)$ be another countable atlas, with pieces
$E'_k$. The sets $B\cap E_i\cap E'_k$ are disjoint with union
$B\cap E_i$, so by \cref{lebesgue-integral:cor:additive}
\[
  \int_{\varphi_i(B\cap E_i)}\rho_{\varphi_i}\,\dd\lambda
  = \sum_k\int_{\varphi_i(B\cap E_i\cap E'_k)}\rho_{\varphi_i}\,
    \dd\lambda
  = \sum_k\int_{\varphi'_k(B\cap E_i\cap E'_k)}\rho_{\varphi'_k}\,
    \dd\lambda .
\]
The second equality is the change of variables theorem
(\cref{change-of-variables:thm:cov}) for the transition map
$\tau = \varphi'_k\circ\varphi_i^{-1}$ and the function
$\rho_{\varphi'_k}$ times the indicator of
$\varphi'_k(B\cap E_i\cap E'_k)$, together with
\eqref{stokes:eq:density}. Summing over $i$, and exchanging the two
sums of non-negative terms, gives the number computed with the second
atlas. If $B\subset U$ for a chart $\varphi$, put this chart first in
an atlas: then $E_1\supset B$ and the sum has one term. If $B$ is the
disjoint union of Borel sets $B_j$, each integral over
$\varphi_i(B\cap E_i)$ is the sum of those over
$\varphi_i(B_j\cap E_i)$, and the sums can again be exchanged. For a
null set every term is the integral over a Lebesgue-null set.
\end{proof}

A density $\rho$ is \emph{integrable} if $\int_M\abs\rho < \infty$, and
then
$\int_B\rho = \int_B\rho^+ - \int_B\rho^-$,
with $\rho^\pm$ the positive and negative parts. The integral is linear
in $\rho$. A function $h$ is integrated \emph{against} a density:
$\int_Mh\,\rho$ is the integral of $h$ for the measure defined by
$\rho$ when $\rho\geq0$. On a Riemannian manifold, the measure of the
volume density is the \emph{Riemannian volume}
$\operatorname{vol}_M$, and $\int_Mh\,\dd\!\operatorname{vol}_M$ is the
integral of the function $h$. No orientation has been used.

\begin{keyidea}
On $\R^n$ a function and a density look the same, because Lebesgue
measure is there to turn one into the other. On a manifold there is no
preferred measure. A probability distribution is a density, or a
function \emph{times a chosen reference density}; changing the
reference changes the function. A Riemannian metric provides a
reference, an inner product on each tangent space being enough to
measure volumes (\cref{multilinear-algebra:tab:structure}).
\end{keyidea}

\section{The integral of a form}\label{stokes:sec:forms}

\begin{definition}[Integral of a form of top degree]
\label{stokes:def:integral}
Let $M$ be oriented and $\omega$ an $n$-form on $M$ with continuous
components. It is \emph{integrable} if the density $\rho^\omega$ of
\cref{stokes:ex:densities} is, and then, for a Borel set $B$,
\[
  \int_B\omega = \int_B\rho^\omega .
\]
\end{definition}

So if $B$ lies in the domain of a positive chart $(U, x)$ and
$\omega = f\,\dd x_1\wedge\dots\wedge\dd x_n$, then
$\int_B\omega = \int_{x(B)}f\,\dd\lambda$, where $f$ is read in the
chart: erase the wedges. In a negative chart a minus sign appears, and
reversing the orientation of $M$ changes the sign of every integral.

\begin{proposition}[Computing integrals]\label{stokes:prop:computing}
Let $M$ be oriented and $\omega$ an integrable $n$-form.
\begin{enumerate}[nosep]
  \item If $F:N\to M$ is a diffeomorphism that preserves orientation,
        then $\int_NF^*\omega = \int_M\omega$.
  \item Let $\psi:\Omega\to M$ be a diffeomorphism from an open subset
        of $\R^n$ onto an open subset of $M$ whose complement is a null
        set, such that $\psi^{-1}$ is a positive chart. If
        $\psi^*\omega = g\,\dd t_1\wedge\dots\wedge\dd t_n$, then
        $\int_M\omega = \int_\Omega g\,\dd\lambda$, with
        \[
          g(t) = \omega\bigl(\psi(t)\bigr)\bigl(\partial_1\psi(t),\dots,
          \partial_n\psi(t)\bigr) .
        \]
\end{enumerate}
\end{proposition}

\begin{proof}
(1) The charts $\varphi\circ F$, for $\varphi$ a positive chart of $M$,
are positive charts of $N$, in which $F^*\omega$ has the same component
as $\omega$ in $\varphi$. (2) Put the chart $\psi^{-1}$ first in a
countable atlas of positive charts. Then $E_1 = \psi(\Omega)$ and the
other pieces lie in a null set. The component of $\omega$ in the chart
$\psi^{-1}$ is $g$.
\end{proof}

A parametrization that misses a null set is thus enough to compute an
integral: the sphere without a meridian, a torus without two circles.

\paragraph{Forms of any degree.} A $k$-form $\omega$ on a manifold $N$
is integrated over an oriented manifold $S$ of dimension $k$ that is
mapped into $N$ by a smooth map $\iota$, for instance the inclusion of
a submanifold:
\[
  \int_S\omega = \int_S\iota^*\omega .
\]
With a positive parametrization $\psi$ of $S$ as in
\cref{stokes:prop:computing}, this is
\[
  \int_S\omega = \int_\Omega\omega\bigl(\psi(t)\bigr)
  \bigl(\partial_1\psi(t),\dots,\partial_k\psi(t)\bigr)\,\dd t .
\]
For a curve $c:(a, b)\to N$ it is
$\int_a^b\omega(c(t))(c'(t))\,\dd t$, the integral used in
\cref{differential-forms:prop:circulation}. The form eats the small
parallelepiped spanned by the velocities of the parametrization, and
the results are added.

\section{Domains with boundary}\label{stokes:sec:boundary}

\begin{definition}[Domain with smooth boundary]\label{stokes:def:domain}
A subset $D$ of $M$ is a \emph{domain with smooth boundary} if every
point $p$ of $M$ has a chart $(U, x)$ around it of one of three kinds:
\begin{enumerate}[nosep,label=(\roman*)]
  \item $U\subset D$;
  \item $U\cap D = \emptyset$;
  \item $x_1(p) = 0$ and $U\cap D = \{q\in U : x_1(q)\leq0\}$.
\end{enumerate}
The \emph{boundary} $\partial D$ is the set of points of the third
kind. Charts of these three kinds are called \emph{adapted} to $D$.
\end{definition}

A point of the third kind has points of $D$ and points outside $D$
arbitrarily close to it, so the three kinds exclude one another. The
set $D$ is closed, and it is Borel, since countably many adapted charts
cover $M$ (\cref{stokes:lem:cover} below). In an adapted chart of kind
(iii), $\partial D\cap U = \{x_1 = 0\}$.

\begin{proposition}[Sublevel sets]\label{stokes:prop:sublevel}
Let $g:M\to\R$ be smooth with $\dd g(p)\neq0$ whenever $g(p) = 0$. Then
$D = \{g\leq0\}$ is a domain with smooth boundary, and
$\partial D = \{g = 0\}$.
\end{proposition}

\begin{proof}
The sets $\{g < 0\}$ and $\{g > 0\}$ are open, which gives charts of
kinds (i) and (ii). Let $g(p) = 0$ and $(U, y)$ be a chart around $p$.
Some $\partial g/\partial y_i(p)$ is not $0$, say for $i = 1$. The map
$(g, y_2,\dots,y_n)$ has an invertible differential at $p$, so it is a
chart on a smaller open set
(\cref{tangent-bundles:thm:inverse}), of kind (iii).
\end{proof}

For instance the closed ball $\{\norm x^2 - 1\leq0\}$ of a Euclidean
space, with boundary the sphere.

\begin{proposition}[The boundary]\label{stokes:prop:boundary}
Let $D$ be a domain with smooth boundary in $M$.
\begin{enumerate}[nosep]
  \item $\partial D$ is a manifold of dimension $n - 1$, with the
        charts $(U\cap\partial D; x_2,\dots,x_n)$ given by the adapted
        charts of kind \textup{(iii)}. The inclusion
        $\iota:\partial D\to M$ is smooth, and $D\iota(p)$ identifies
        $T_p\partial D$ with $\ker\dd x_1(p)\subset T_pM$.
  \item For $p\in\partial D$, the sign of $\dd x_1(p)v$ is the same in
        all adapted charts around $p$. A vector $v\in T_pM$ is
        \emph{outward} if $\dd x_1(p)v > 0$.
\end{enumerate}
\end{proposition}

\begin{proof}
Let $x$ and $x'$ be two adapted charts of kind (iii) around
$p\in\partial D$, and $\tau$ the transition map. It maps $\{x_1 = 0\}$
into $\{x'_1 = 0\}$ and $\{x_1\leq0\}$ into $\{x'_1\leq0\}$. Hence, at
the points with $x_1 = 0$,
\[
  \frac{\partial x'_1}{\partial x_j} = 0\quad(j\geq2) ,
  \qquad
  \frac{\partial x'_1}{\partial x_1}\geq0 ,
\]
and the second number is not $0$, since the Jacobian matrix $J$ is
invertible. So $\dd x'_1(p) = c\,\dd x_1(p)$ with $c > 0$, which is (2).
The first row of $J$ is $(c, 0,\dots,0)$, so
\begin{equation}\label{stokes:eq:block}
  \det J = c\,\det\tilde J ,
\end{equation}
where $\tilde J$, the Jacobian matrix of
$(x_2,\dots,x_n)\mapsto(x'_2,\dots,x'_n)$ on $\{x_1 = 0\}$, is
invertible. These restricted transition maps are smooth, with smooth
inverses, between open subsets of $\R^{n-1}$: the charts of (1) form an
atlas. Two points of $\partial D$ are separated by open sets of $M$,
and countably many of these charts suffice
(\cref{stokes:lem:cover} below). In the charts $x$ and
$(x_2,\dots,x_n)$, $\iota$ reads $y\mapsto(0, y)$.
\end{proof}

\begin{definition}[Boundary orientation]\label{stokes:def:bdry-orientation}
Let $M$ be oriented and $D$ a domain with smooth boundary. For
$n\geq2$, a basis $(v_2,\dots,v_n)$ of $T_p\partial D$ is
\emph{positive} if $(\nu, v_2,\dots,v_n)$ is a positive basis of $T_pM$
for an outward vector $\nu$. For $n = 1$, a point $p$ of $\partial D$
has the sign $+1$ if the outward vectors at $p$ are positive, and $-1$
otherwise.
\end{definition}

In words: \emph{outward first}. The choice of $\nu$ does not matter:
two outward vectors differ by a positive factor and a vector tangent to
$\partial D$, which does not change the sign of
$\mu(\nu, v_2,\dots,v_n)$ for an $n$-form $\mu$. This is an orientation
of $\partial D$. Indeed, around $p\in\partial D$ take an adapted chart
of kind (iii); if it is not positive for $M$, replace $x_2$ by $-x_2$.
Now $\partial_1$ is outward and $(\partial_1,\dots,\partial_n)$ is
positive, so $(\partial_2,\dots,\partial_n)$ is positive on
$\partial D\cap U$: the chart $(x_2,\dots,x_n)$ of $\partial D$ is
positive. For the ball this is the orientation of the sphere of
\cref{stokes:ex:orientations}. For a domain of the oriented plane, the
boundary is traversed with the domain on its left
(\cref{stokes:fig:domain}).

\begin{figure}[ht]
\centering
\begin{tikzpicture}[>=Stealth, scale=1.0]
  \begin{scope}
    \fill[bookintuition!16] plot[smooth cycle, tension=0.8]
      coordinates {(0,0) (1.6,-0.7) (3.4,-0.2) (3.9,1.4) (2.6,2.6)
      (0.8,2.3) (-0.3,1.2)};
    \draw[bookaccent, very thick,
      decoration={markings,
        mark=at position 0.12 with {\arrow{Stealth}},
        mark=at position 0.45 with {\arrow{Stealth}},
        mark=at position 0.8 with {\arrow{Stealth}}},
      postaction={decorate}] plot[smooth cycle, tension=0.8]
      coordinates {(0,0) (1.6,-0.7) (3.4,-0.2) (3.9,1.4) (2.6,2.6)
      (0.8,2.3) (-0.3,1.2)};
    \node at (1.7,1.0) {$D$};
    \node[bookaccent] at (0.2,2.55) {$\partial D$};
    \fill (3.9,1.4) circle (1.6pt) node[below left=-1pt] {$p$};
    \draw[->, thick, bookpitfall] (3.9,1.4) -- ++(0.95,0.2)
      node[right] {$\nu$};
    \draw[->, thick] (3.9,1.4) -- ++(-0.2,0.95) node[above] {$v_2$};
    \draw[bookdef, dashed] (3.9,1.4) circle (0.75);
    \node[bookdef] at (4.55,0.45) {$U$};
  \end{scope}
  \draw[->, thick, bookdef] (5.2,1.9) to[bend left=18]
    node[above] {$x$} (7.0,1.9);
  \begin{scope}[shift={(9.2,1.2)}]
    \fill[bookintuition!16] (-1.6,-1.3) rectangle (0,1.3);
    \draw[bookaccent, very thick, ->] (0,-1.3) -- (0,1.3);
    \draw[->] (-1.9,0) -- (1.5,0) node[right] {$x_1$};
    \node[bookaccent, above] at (0,1.3) {$x_2$};
    \node at (-0.85,-0.75) {$x_1\leq0$};
    \fill (0,0) circle (1.6pt);
    \draw[->, thick, bookpitfall] (0,0) -- (0.95,0)
      node[above] {$\partial_1$};
    \draw[->, thick] (0,0) -- (0,0.95) node[left=1pt] {$\partial_2$};
  \end{scope}
\end{tikzpicture}
\caption{A domain with smooth boundary in an oriented surface, and an
adapted chart at a boundary point: the domain becomes the half-plane
$x_1\leq0$. The vector $\partial_1$ is outward. The boundary is oriented
so that an outward vector followed by a positive basis of the boundary
is a positive basis of the surface: outward first.}
\label{stokes:fig:domain}
\end{figure}

To prove a statement about a form on a domain, one cuts the form into
pieces, each of which lives in one adapted chart. Three lemmas prepare
this. A subset $K$ of $M$ is \emph{compact} if every sequence of points
of $K$ has a subsequence that converges to a point of $K$. The
\emph{support} of a form is the smallest closed set outside which it
vanishes.

\begin{lemma}[Smooth steps]\label{stokes:lem:step}
\begin{enumerate}[nosep]
  \item The function $h(t) = e^{-1/t}$ for $t > 0$, $h(t) = 0$ for
        $t\leq0$, is smooth on $\R$.
  \item $\theta(t) = h(1 - t)/\bigl(h(1 - t) + h(t - \frac12)\bigr)$ is
        smooth, with values in $[0, 1]$, equal to $1$ for
        $t\leq\frac12$ and to $0$ for $t\geq1$.
  \item $\beta(x) = h\bigl(1 - \sum_ix_i^2\bigr)$ is smooth on $\R^n$,
        positive on the open unit ball and $0$ outside.
\end{enumerate}
\end{lemma}

\begin{proof}
(1) For $t > 0$, by induction, $h^{(k)}(t) = P_k(1/t)\,e^{-1/t}$ with
$P_k$ a polynomial, since the derivative of $P(1/t)\,e^{-1/t}$ is
$\bigl(P(1/t) - P'(1/t)\bigr)t^{-2}e^{-1/t}$. As $s^me^{-s}\to0$ when
$s\to\infty$, both $h^{(k)}(t)$ and $h^{(k)}(t)/t$ tend to $0$ as
$t\to0^+$. So, by induction on $k$, $h^{(k)}$ exists at $0$, equals
$0$ there, and is continuous. (2) The denominator is positive: its
first term is when $t < 1$, its second when $t > \frac12$. (3) is a
composition of smooth maps.
\end{proof}

\begin{lemma}[Extracting covers]\label{stokes:lem:cover}
Let $\mathcal W$ be a family of open subsets of $M$ whose union
contains a set $A$. Then countably many members of $\mathcal W$ cover
$A$, and finitely many if $A$ is compact.
\end{lemma}

\begin{proof}
Take a countable atlas $(U_i, \varphi_i)$ and let $\mathcal B$ be the
countable family of the sets $\varphi_i^{-1}(B)$, for $B$ an open ball
with rational centre and radius contained in $\varphi_i(U_i)$. Every
open set $W$ is a union of members of $\mathcal B$: if $p\in W\cap U_i$,
the open set $\varphi_i(W\cap U_i)$ contains such a ball around
$\varphi_i(p)$. The members of $\mathcal B$ that lie in some member of
$\mathcal W$ therefore cover $A$; choosing for each of them one member
of $\mathcal W$ that contains it gives a countable subfamily
$W_1, W_2,\dots$ covering $A$. Let $A$ be compact and suppose that no
$W_1\cup\dots\cup W_N$ contains $A$. Choose
$a_N\in A\setminus(W_1\cup\dots\cup W_N)$. A subsequence converges to
some $a\in A$, which lies in some $W_m$. Since $W_m$ is open, it
contains $a_N$ for infinitely many $N\geq m$: a contradiction.
\end{proof}

\begin{lemma}[Partition of unity on a compact set]
\label{stokes:lem:pou}
Let $K\subset M$ be compact and $\mathcal W$ a family of open sets
covering $K$. There are smooth functions
$\rho_1,\dots,\rho_N:M\to[0, 1]$, each with compact support contained
in a member of $\mathcal W$, such that
$\rho_1 + \dots + \rho_N = 1$ on $K$.
\end{lemma}

\begin{proof}
For $p\in K$ choose a member $W$ of $\mathcal W$ and a chart
$(U, \varphi)$ around $p$, and $r > 0$ such that the closed ball of
centre $\varphi(p)$ and radius $r$ lies in $\varphi(U\cap W)$. Define
$b_p = \beta\bigl((\varphi - \varphi(p))/r\bigr)$ on $U$ and $b_p = 0$
elsewhere. The preimage $C$ of the closed ball is compact, hence
closed: a limit of points of $C$ is the limit of a subsequence in $C$,
and limits are unique in $M$ by the Hausdorff property. So $b_p$ is
smooth on $M$, being smooth on $U$ and zero on the open set
$M\setminus C$; its support is $C\subset W$, and $b_p > 0$ on an open
set $W_p\ni p$. By \cref{stokes:lem:cover} finitely many $W_p$ cover
$K$; let $b_1,\dots,b_N$ be the corresponding functions and
$\sigma = b_1 + \dots + b_N$. It is positive on $K$, and attains a
smallest value $c > 0$ there, as in
\cref{norms-banach:thm:extreme}. With $\theta$ from
\cref{stokes:lem:step}, put
\[
  \rho_j = \frac{b_j}{\sigma + \theta(\sigma/c)} .
\]
The denominator is positive everywhere, since
$\theta(\sigma/c) = 1$ where $\sigma\leq c/2$. On $K$, $\sigma\geq c$,
so $\theta(\sigma/c) = 0$ and $\sum_j\rho_j = 1$.
\end{proof}

In particular a continuous $n$-form $\omega$ with compact support on an
oriented manifold is integrable: with $K$ its support and $\mathcal W$
the domains of the positive charts, $\omega = \sum_j\rho_j\omega$, and
each $\rho_j\omega$ has a bounded component that vanishes outside a
bounded set in one chart.

\section{Stokes' theorem}\label{stokes:sec:stokes}

\begin{theorem}[Stokes]\label{stokes:thm:stokes}
Let $M$ be an oriented manifold of dimension $n$, $D\subset M$ a domain
with smooth boundary, and $\omega$ a smooth $(n - 1)$-form on $M$.
Assume that $\omega$ has compact support, or that $D$ is compact. Then
\[
  \int_D\dd\omega = \int_{\partial D}\omega ,
\]
where $\partial D$ has the boundary orientation and the right-hand side
is the integral of the restriction $\iota^*\omega$. For $n = 1$ the
right-hand side is the sum of the values of the function $\omega$ at
the points of $\partial D$, with their signs.
\end{theorem}

\begin{proof}
\emph{Reduction to one chart.} Let $K$ be the support of $\omega$, or
$D$ if $D$ is compact, and $\mathcal W$ the family of the domains of
the adapted charts that are positive or, in dimension $1$, negative.
They cover $M$: an adapted chart has a constant sign on a smaller open
set around a given point, and if $n\geq2$ a negative one becomes
positive when $x_2$ is replaced by $-x_2$. With $\rho_j$ from
\cref{stokes:lem:pou}, let
$\omega' = \sum_j\rho_j\omega$. In the first case $\omega' = \omega$.
In the second, $\omega' = \omega$ on $D$; then
$\dd\omega' = \dd\omega$ at the points of $D$ that have a chart of kind
(i), and by continuity at the points of $\partial D$, which are limits
of such points. Both sides of the formula are thus the same for
$\omega$ and $\omega'$, and both are linear. So we may assume that
$\omega$ has compact support in the domain $U$ of an adapted chart
$x$, positive if $n\geq2$. In this chart
\[
  \omega = \sum_{i=1}^nf_i\,
  \dd x_1\wedge\dots\wedge\widehat{\dd x_i}\wedge\dots\wedge\dd x_n ,
  \qquad
  \dd\omega = \sum_{i=1}^n(-1)^{i-1}\,\partial_if_i\;
  \dd x_1\wedge\dots\wedge\dd x_n ,
\]
where the $f_i$, read in the chart, are smooth functions on $\R^n$ that
vanish outside a bounded set.

\emph{Kind (ii).} Both sides are $0$.

\emph{Kind (i).} The right-hand side is $0$, since
$U\cap\partial D = \emptyset$. The left-hand side is
$\sum_i(-1)^{i-1}\int_{\R^n}\partial_if_i\,\dd\lambda$. By Fubini's
theorem (\cref{fubini:thm:fubini}) integrate first in $x_i$: the inner
integral $\int_\R\partial_if_i\,\dd x_i$ is $0$ by the fundamental
theorem of calculus, because $f_i$ vanishes for $\abs{x_i}$ large.

\emph{Kind (iii), $n\geq2$.} The left-hand side is
$\sum_i(-1)^{i-1}\int_{\{x_1\leq0\}}\partial_if_i\,\dd\lambda$. For
$i\geq2$ the inner integral in $x_i$ over $\R$ is $0$ as before. For
$i = 1$,
\[
  \int_{-\infty}^0\partial_1f_1(x_1, y)\,\dd x_1 = f_1(0, y) ,
  \qquad y = (x_2,\dots,x_n) ,
\]
so the left-hand side is $\int_{\R^{n-1}}f_1(0, y)\,\dd y$. On the
other side, $x_1\circ\iota = 0$ gives $\iota^*\dd x_1 = 0$, so
\[
  \iota^*\omega = f_1(0, \cdot)\;\dd x_2\wedge\dots\wedge\dd x_n ,
\]
and $(x_2,\dots,x_n)$ is a positive chart of $\partial D$. The
right-hand side is the same integral.

\emph{Kind (iii), $n = 1$.} Here $\omega = f$ is a function and
$U\cap D = \{x\leq0\}$. If the chart is positive, the left-hand side is
$\int_{-\infty}^0f'(x)\,\dd x = f(0)$, and the point has the sign $+1$
because $\partial_x$ is outward and positive. If the chart is negative,
the left-hand side is $-f(0)$ and the sign of the point is $-1$.
\end{proof}

\begin{intuition}
In a chart the theorem is the fundamental theorem of calculus in the
direction that crosses the boundary, and Fubini's theorem in the
others. The partition of unity does the bookkeeping that
\cref{ch:differential-forms} described with small parallelograms: the
contributions of the pieces of the boundary that lie inside the domain
cancel, and those on $\partial D$ remain.
\end{intuition}

\begin{corollary}\label{stokes:cor:consequences}
Let $M$ be oriented, of dimension $n$.
\begin{enumerate}[nosep]
  \item If $\eta$ is a smooth $(n - 1)$-form with compact support, then
        $\int_M\dd\eta = 0$.
  \item \emph{Integration by parts.} For a function $f$ and an
        $(n - 1)$-form $\eta$, under the hypotheses of the theorem,
        \[
          \int_D\dd f\wedge\eta
          = \int_{\partial D}f\,\eta - \int_Df\,\dd\eta .
        \]
  \item If $\omega$ is closed on $M$ and $D$ is compact, then
        $\int_{\partial D}\omega = 0$.
\end{enumerate}
\end{corollary}

\begin{proof}
(1) is the theorem for $D = M$, whose boundary is empty. (2) is the
theorem for $f\eta$, with
$\dd(f\eta) = \dd f\wedge\eta + f\,\dd\eta$. (3) is the theorem.
\end{proof}

\begin{example}[The plane]\label{stokes:ex:green}
Let $D$ be a compact domain with smooth boundary in $\R^2$. For
$\omega = P\,\dd x + Q\,\dd y$, the theorem reads
\[
  \int_{\partial D}P\,\dd x + Q\,\dd y
  = \int_D\Bigl(\frac{\partial Q}{\partial x}
    - \frac{\partial P}{\partial y}\Bigr)\,\dd x\,\dd y ,
\]
Green's formula. With $\omega = \frac12(x\,\dd y - y\,\dd x)$, whose
exterior derivative is $\dd x\wedge\dd y$, the area of $D$ is the
integral of $\omega$ along its boundary.
\end{example}

\begin{example}[Sphere and ball]\label{stokes:ex:sphere}
In $\R^3$ let $\mu = \dd x\wedge\dd y\wedge\dd z$ and
$\sigma = i_E\,\mu = x\,\dd y\wedge\dd z + y\,\dd z\wedge\dd x
+ z\,\dd x\wedge\dd y$, with $E$ the radial field. Then
$\dd\sigma = 3\mu$, and on the unit ball $B$
\[
  \int_{S^2}\sigma = \int_B3\mu = 3\,\lambda(B) = 4\pi .
\]
Directly: $\psi(\theta, \phi)
= (\sin\theta\cos\phi, \sin\theta\sin\phi, \cos\theta)$ on
$(0, \pi)\times(0, 2\pi)$ misses a null set, and
$\psi^*\sigma = \mu(\psi, \partial_\theta\psi, \partial_\phi\psi)\,
\dd\theta\wedge\dd\phi = \sin\theta\,\dd\theta\wedge\dd\phi$. The
coefficient is positive, so the chart is positive for the outward
orientation, and $\int_0^{2\pi}\!\int_0^\pi\sin\theta\,\dd\theta\,
\dd\phi = 4\pi$. On unit vectors tangent to the sphere, $\sigma$ is the
area: the two computations give the area of the sphere and three times
the volume of the ball, as in \cref{change-of-variables:cor:ball}.
\end{example}

\begin{example}[Winding]\label{stokes:ex:winding}
Let $\omega$ be the angle form of \cref{differential-forms:ex:angle}
and $D$ a compact domain with smooth boundary in $\R^2$, with
$0\notin\partial D$. If $0\notin D$, then $D$ is a domain in the
punctured plane, where $\omega$ is closed, and
$\int_{\partial D}\omega = 0$. If $0\in D$, remove from $D$ a small
open disc centred at $0$. The result is a domain in the punctured
plane whose boundary is $\partial D$ together with a circle, on which
the outward vectors of the domain point towards $0$: the circle is
traversed clockwise (\cref{stokes:fig:winding}). Hence
$\int_{\partial D}\omega - 2\pi = 0$. So
\[
  \frac1{2\pi}\int_{\partial D}\omega =
  \begin{cases}1 & \text{if }0\in D ,\\ 0 & \text{if }0\notin D .
  \end{cases}
\]
The integral of a closed form over a boundary does not see the shape of
the boundary, only the holes that it encloses.
\end{example}

\begin{figure}[ht]
\centering
\begin{tikzpicture}[>=Stealth, scale=1.25]
  \fill[bookintuition!16] plot[smooth cycle, tension=0.8]
    coordinates {(-2.3,-0.4) (-0.6,-1.5) (1.8,-1.2) (2.6,0.3)
    (1.2,1.6) (-1.0,1.4)};
  \draw[bookaccent, very thick,
    decoration={markings,
      mark=at position 0.1 with {\arrow{Stealth}},
      mark=at position 0.4 with {\arrow{Stealth}},
      mark=at position 0.75 with {\arrow{Stealth}}},
    postaction={decorate}] plot[smooth cycle, tension=0.8]
    coordinates {(-2.3,-0.4) (-0.6,-1.5) (1.8,-1.2) (2.6,0.3)
    (1.2,1.6) (-1.0,1.4)};
  \fill[white] (0,0) circle (0.55);
  \draw[bookpitfall, very thick,
    decoration={markings, mark=at position 0.3 with {\arrow{Stealth}}},
    postaction={decorate}]
    (0.55,0) arc (0:-360:0.55);
  \fill (0,0) circle (1.4pt) node[below=1pt] {\small$0$};
  \draw[->, thick] (0.39,0.39) -- (0.1,0.1);
  \node at (0.85,0.75) {\small outward};
  \node[bookaccent] at (2.75,1.2) {$\partial D$};
  \node at (1.45,-0.62) {$D$ minus a disc};
\end{tikzpicture}
\caption{A domain of the plane from which a small disc around the
origin has been removed. The boundary of the new domain is the old
boundary, traversed with the domain on its left, and the small circle,
also with the domain on its left: clockwise.}
\label{stokes:fig:winding}
\end{figure}

\begin{remark}[Corners]\label{stokes:rem:corners}
A square or a simplex is not a domain with smooth boundary. The theorem
holds for them as well, the integral over the boundary being taken over
its smooth part: Loomis and Sternberg prove it for domains with
``almost regular'' boundary and Hubbard and Hubbard for
``pieces-with-corners''. For a rectangle it follows from Fubini's
theorem alone (\cref{stokes:exr:5}).
\end{remark}

\section{A discrete model: graphs}\label{stokes:sec:graphs}

A graph is a one-dimensional object made of points and segments, and
the calculus of this chapter exists on it in a finite form, where the
theorems become statements of linear algebra.

Let a graph have a finite set $N$ of nodes and a finite set $E$ of
edges, each edge $e$ being directed from a node $e^-$ to a node $e^+$.
The directions are a bookkeeping device, like an orientation. A
\emph{$0$-cochain} is a function on $N$ and a \emph{$1$-cochain} a
function on $E$; they form the spaces $\R^N$ and $\R^E$. The
\emph{differential} is
\[
  \dd:\R^N\to\R^E , \qquad (\dd f)(e) = f(e^+) - f(e^-) .
\]
Its matrix, with one row per edge and one column per node, has $-1$ at
the tail and $+1$ at the head of each edge: the \emph{incidence matrix}
of the graph. A \emph{path} is a sequence of edges in which each one
starts, after a possible reversal, where the previous one ends; the
integral of a $1$-cochain $\alpha$ along it is the sum of its values on
these edges, with a minus sign for the reversed ones. A path is
\emph{closed} if it ends where it starts. Two nodes are in the same
\emph{component} if a path joins them.

\begin{proposition}[Calculus on a graph]\label{stokes:prop:graph}
Let the graph have $n$ nodes, $m$ edges and $c$ components.
\begin{enumerate}[nosep]
  \item The integral of $\dd f$ along a path from $a$ to $b$ is
        $f(b) - f(a)$.
  \item $\dd f = 0$ if and only if $f$ is constant on each component.
        So $\dd$ has rank $n - c$.
  \item A $1$-cochain is of the form $\dd f$ if and only if its
        integral along every closed path is $0$.
  \item The dual operator $\dual\dd$ sends a covector $y$ on $\R^E$, a
        row of numbers $y_e$, to the covector on $\R^N$ with components
        \[
          (\dual\dd y)_a = \sum_{e^+ = a}y_e - \sum_{e^- = a}y_e .
        \]
        Its kernel is the annihilator of $\im\dd$ and has dimension
        $m - n + c$.
\end{enumerate}
\end{proposition}

\begin{proof}
(1) The sum telescopes. (2) By (1), $\dd f = 0$ if and only if $f$ takes
the same value at any two nodes joined by a path; the kernel has
dimension $c$, and the rank follows from the rank--nullity theorem.
(3) If $\alpha = \dd f$, use (1). Conversely choose a node $a_0$ in
each component and let $f(a)$ be the integral of $\alpha$ along a path
from the chosen node of the component of $a$ to $a$. Two such paths,
one followed by the other reversed, make a closed path, so $f(a)$ is
well defined; and for an edge $e$, extending a path to $e^-$ by $e$
gives $f(e^+) = f(e^-) + \alpha(e)$. (4) By definition
$(\dual\dd y)(f) = y(\dd f) = \sum_ey_e\bigl(f(e^+) - f(e^-)\bigr)$;
collect the coefficient of each $f(a)$. The rest is
\cref{duality:thm:four} and \cref{duality:prop:annihilator}.
\end{proof}

\begin{keyidea}
Read a covector $y$ on $\R^E$ as a \emph{chain}: a combination of
edges with coefficients, over which $1$-cochains are integrated,
$\int_y\alpha = y(\alpha)$. A path is the chain with coefficients
$\pm1$ on its edges. Then $\dual\dd y$ is the \emph{boundary} of the
chain, a combination of nodes: for one edge, its head minus its tail.
The definition of the dual operator,
\[
  \int_y\dd f = (\dual\dd y)(f) = \int_{\partial y}f ,
\]
is Stokes' theorem. On a graph the theorem is not proved: the boundary
is \emph{defined} as the dual of the differential. A chain without
boundary is a \emph{cycle}; (3) and (4) say that the exact cochains are
those that integrate to $0$ on the cycles, and that there are
$m - n + c$ independent cycles.
\end{keyidea}

\begin{figure}[ht]
\centering
\begin{tikzpicture}[>=Stealth, scale=1.0,
    nd/.style={circle, draw, thick, fill=bookintuition!14,
      inner sep=2pt, minimum size=17pt}]
  \node[nd] (1) at (0,0) {$1$};
  \node[nd] (2) at (2.6,1.5) {$2$};
  \node[nd] (3) at (2.6,-1.5) {$3$};
  \node[nd] (4) at (5.2,0) {$4$};
  \draw[->, thick] (1) -- node[above left=-2pt] {$e_1$} (2);
  \draw[->, thick] (1) -- node[below left=-2pt] {$e_2$} (3);
  \draw[->, thick] (2) -- node[right=-1pt] {$e_3$} (3);
  \draw[->, thick] (2) -- node[above right=-2pt] {$e_4$} (4);
  \draw[->, thick] (3) -- node[below right=-2pt] {$e_5$} (4);
  \draw[->, bookpitfall, thick] (1.3,0) ++(50:0.38)
    arc (50:-260:0.38);
  \draw[->, bookpitfall, thick] (3.9,0) ++(130:0.38)
    arc (130:440:0.38);
  \node at (9.3,0) {$\dd = \begin{bmatrix}
    -1 & 1 & 0 & 0\\
    -1 & 0 & 1 & 0\\
    0 & -1 & 1 & 0\\
    0 & -1 & 0 & 1\\
    0 & 0 & -1 & 1
  \end{bmatrix}$};
\end{tikzpicture}
\caption{A graph with $n = 4$ nodes, $m = 5$ edges and $c = 1$
component, and its incidence matrix. The rank is $n - c = 3$. There are
$m - n + c = 2$ independent cycles, for instance
$e_1 + e_3 - e_2$ and $e_3 + e_5 - e_4$, drawn as the two loops.}
\label{stokes:fig:graph}
\end{figure}

With the dot products of $\R^N$ and $\R^E$, the adjoint of $\dd$ has
the matrix $[\dd]^\top$, and
\[
  L = \adj\dd\,\dd , \qquad
  (Lf)(a) = \sum_{e\text{ at }a}\bigl(f(a) - f(\text{other end})\bigr) ,
  \qquad
  \inner f{Lf} = \sum_e\bigl(\dd f(e)\bigr)^2 ,
\]
is the \emph{Laplacian} of the graph, the degree matrix minus the
adjacency matrix when no edge is repeated. It does not depend on the
directions of the edges. Its kernel is the kernel of $\dd$: the number
of components is the multiplicity of the eigenvalue $0$. Assigning to
each electrical wire its current, a covector $y$, the equation
$\dual\dd y = 0$ is Kirchhoff's current law, and (3) is his voltage
law: the voltage drops come from a potential.

\begin{mlthread}[Machine learning: learning on graphs]
\begin{itemize}[nosep]
  \item The quadratic form $\inner f{Lf}$ measures how much a function
        on the nodes varies along the edges. Penalizing it spreads
        labels from the labelled nodes to the others, and one step of
        gradient descent on it, $f\leftarrow f - \eta Lf$, replaces the
        value at a node by an average with its neighbours: the
        propagation step of graph neural networks.
  \item The eigenvectors of $L$ for its smallest eigenvalues are the
        slowest varying functions on the graph, minimizers of the
        Rayleigh quotient (\cref{ch:spectral-svd}); they are used to
        cut a graph into clusters and as coordinates for its nodes.
  \item These subjects are developed in Part~XI of the companion
        \emph{Notes on Machine Learning} (Graph Laplacians and Spectral
        Methods, and the chapters that follow it). What this section
        adds is the factorization $L = \adj\dd\,\dd$ and its meaning.
\end{itemize}
\end{mlthread}

\section{Integration by parts in machine learning}
\label{stokes:sec:ml}

On $\R^n$, let $\mu = \dd x_1\wedge\dots\wedge\dd x_n$ and, for a
vector field $X = \sum_iX_i\,\partial_i$, let
$\operatorname{div}X = \sum_i\partial_iX_i$.

\begin{proposition}[Integration by parts on $\R^n$]
\label{stokes:prop:parts}
Let $f$ be a smooth function and $X$ a smooth vector field on $\R^n$
such that $fX$ has compact support. Then
\[
  \int_{\R^n}\bigl(\dd f(X) + f\operatorname{div}X\bigr)\,\dd\lambda
  = 0 .
\]
\end{proposition}

\begin{proof}
Let $\eta = f\,i_X\mu$, an $(n - 1)$-form with compact support. By (2)
of \cref{differential-forms:prop:interior},
$i_X\mu = \sum_i(-1)^{i-1}X_i\,\dd x_1\wedge\dots\wedge
\widehat{\dd x_i}\wedge\dots\wedge\dd x_n$, so
$\dd(i_X\mu) = (\operatorname{div}X)\,\mu$. By (1) of the same
proposition, $0 = i_X(\dd f\wedge\mu)
= \dd f(X)\,\mu - \dd f\wedge i_X\mu$. Hence
$\dd\eta = \dd f\wedge i_X\mu + f\,\dd(i_X\mu)
= \bigl(\dd f(X) + f\operatorname{div}X\bigr)\mu$, and
$\int\dd\eta = 0$ by \cref{stokes:cor:consequences}.
\end{proof}

\begin{mlthread}[Machine learning: Stein's identity and score matching]
Let $p > 0$ be a smooth probability density on $\R^n$, with score
$s_p = \nabla\ln p$, and $X$ a smooth vector field such that $pX$ has
compact support, or decreases fast enough for the argument to pass to
the limit over large balls.
\begin{itemize}[nosep]
  \item With $f = p$ the proposition gives \emph{Stein's identity},
        \[
          \E_p\bigl[\inner{s_p}X + \operatorname{div}X\bigr] = 0 ,
        \]
        since $\dd p(X) = p\,\inner{s_p}X$. Every vector field gives a
        function whose mean under $p$ is known to be $0$, computed from
        the score alone: the normalizing constant of $p$ is not needed.
  \item \emph{Score matching.} To fit a model $s$ to the unknown score
        one would minimize $\frac12\E_p\norm{s - s_p}^2$, which
        contains $s_p$. Expanding the square and applying the identity
        to $X = s$,
        \[
          \tfrac12\E_p\norm{s - s_p}^2
          = \E_p\bigl[\tfrac12\norm s^2 + \operatorname{div}s\bigr]
          + \tfrac12\E_p\norm{s_p}^2 .
        \]
        The last term does not depend on $s$, and the first is a mean
        over the data of quantities computed from the model.
  \item The divergence costs $n$ backward passes. It is replaced by
        the random estimate $u^\top Ju$, with $J$ the Jacobian matrix
        of $s$ and $u$ a standard Gaussian vector, whose mean is
        $\tr J$, or avoided by adding noise to the data.
\end{itemize}
\end{mlthread}

\begin{mlthread}[Machine learning: probability on curved spaces]
Directions, rotations and other constrained quantities live on
manifolds.
\begin{itemize}[nosep]
  \item A probability distribution on a Riemannian manifold is
        described by a function $q$ with
        $\int_Mq\,\dd\!\operatorname{vol}_M = 1$. In a chart its
        density for Lebesgue measure is
        $(q\circ\varphi^{-1})\sqrt{\det G}$: the factor
        $\sqrt{\det G}$ belongs to the volume, not to $q$. Forgetting
        it gives a different distribution in each chart.
  \item On the sphere, the uniform distribution is the normalized
        area, and it is the push-forward of a standard Gaussian by
        $x\mapsto x/\norm x$, because the Gaussian density depends
        only on $\norm x$ (\cref{change-of-variables:thm:polar}). The
        von
        Mises--Fisher distribution has the function
        $q(x)\propto e^{\kappa\inner mx}$ with respect to the area.
  \item A flow $g$ between Riemannian manifolds changes densities by
        the ratio of the volume densities,
        $\sqrt{\det(J^\top G'J)}/\sqrt{\det G}$, which replaces
        $\abs{\det Dg}$ in \cref{change-of-variables:thm:pushforward}.
  \item None of this needs an orientation: probabilities are
        densities, not forms. Orientation enters with Stokes' theorem.
\end{itemize}
\end{mlthread}

\section{Where the chapter leads}\label{stokes:sec:outlook}

\Cref{ch:vector-calculus} reads Stokes' theorem in a Euclidean space,
where $1$-forms and $(n - 1)$-forms become vector fields and the
theorem becomes the divergence theorem and the classical theorem of
Stokes on circulation and curl. \Cref{ch:odes-flows} lets domains move
with a flow: the derivative of $\int_{D_t}\omega$ involves the form
$\dd(i_X\omega) + i_X\dd\omega$ announced in
\cref{ch:differential-forms}, and for a density it gives the equation
that governs continuous normalizing flows.

\paragraph{Sources.} Loomis and Sternberg, \emph{Advanced Calculus},
build partitions of unity in \S10.2 (Theorem~2.1), define densities by
the transition law \eqref{stokes:eq:density} and integrate them in
\S10.3 (Definition~3.1, Theorem~3.1), and define the volume density of
a Riemann metric in \S10.4. Their \S11.2 defines oriented atlases and
positive charts, identifies $n$-forms with densities on an oriented
manifold (Proposition~2.1) and integrates them (Theorem~2.1); \S11.4
orients the boundary of a domain with regular boundary and proves
Stokes' theorem for it (Theorem~4.1). Here the integral of a density
is built as a measure, from a countable atlas, instead of a linear
functional on densities of compact support. Hubbard and Hubbard
integrate forms over parametrized domains (Definition~6.2.1) and
oriented manifolds (Definition~6.4.12), define orientation as a
continuous choice of orientations of the tangent spaces
(Definition~6.3.3), orient the boundary by an outward vector placed
first (Definition~6.6.21), and prove the generalized Stokes theorem for
pieces-with-boundary (Theorem~6.10.2). The reading of Stokes' theorem
on a graph follows the standard language of chains and cochains. On
Stein's identity and score matching: C.~Stein, A bound for the error in
the normal approximation to the distribution of a sum of dependent
random variables, \emph{Proc.\ Sixth Berkeley Symp.}, 1972;
A.~Hyv\"arinen, Estimation of non-normalized statistical models by
score matching, \emph{J.\ Mach.\ Learn.\ Res.}, 2005; Y.~Song et al.,
Sliced score matching, UAI, 2019. On the sphere:
T.~Davidson et al., Hyperspherical variational auto-encoders, UAI,
2018.

\paragraph{Notebook.} The numerical claims of this chapter are checked
with PyTorch in \notebook{23-stokes.ipynb}.

\begin{chaptersummary}
\begin{itemize}[nosep]
  \item An orientation is a continuous choice of positive bases; an
        atlas with positive Jacobian determinants, or an $n$-form
        without zero, gives one. The M\"obius band has none.
  \item A density transforms with $\abs{\det}$ and defines a measure on
        the manifold, without any orientation. A Riemannian metric
        gives the volume density $\sqrt{\det G}$.
  \item On an oriented manifold an $n$-form is a density. In a positive
        chart, $\int f\,\dd x_1\wedge\dots\wedge\dd x_n
        = \int f\,\dd\lambda$. A $k$-form is integrated over an
        oriented $k$-dimensional manifold by pulling it back.
  \item A domain with smooth boundary looks like a half-space in
        adapted charts. Its boundary is oriented by the rule ``outward
        first''.
  \item Stokes: $\int_D\dd\omega = \int_{\partial D}\omega$. The proof
        is a partition of unity, Fubini's theorem and the fundamental
        theorem of calculus.
  \item On a graph, $\dd$ is the incidence matrix, the boundary is its
        dual, and the Laplacian is $\adj\dd\,\dd$.
  \item Integration by parts on $\R^n$ gives Stein's identity and makes
        the score matching objective computable.
\end{itemize}
\end{chaptersummary}

\section*{Exercises}
\addcontentsline{toc}{section}{Exercises}

\begin{exercise}\label{stokes:exr:1}
The sphere $S^2\subset\R^3$ without its north pole, and without its
south pole, has the stereographic charts
$\varphi_N(x) = (x_1, x_2)/(1 - x_3)$ and
$\varphi_S(x) = (x_1, x_2)/(1 + x_3)$.
\begin{enumerate}[nosep]
  \item Show that the transition map is $u\mapsto u/\norm u^2$ and
        that its Jacobian determinant is $-1/\norm u^4$.
  \item Deduce an atlas of two charts that defines an orientation of
        the sphere.
\end{enumerate}
\solutionlink{stokes:exr:1}
\end{exercise}

\begin{exercise}\label{stokes:exr:2}
Let $S$ be the graph $z = x^2 + y^2$ over the disc $x^2 + y^2 < 1$,
oriented so that $(x, y)$ is a positive chart, and
$\omega = z\,\dd x\wedge\dd y - x\,\dd y\wedge\dd z$. Compute
$\int_S\omega$. What is the result for the other orientation?
\solutionlink{stokes:exr:2}
\end{exercise}

\begin{exercise}\label{stokes:exr:3}
\begin{enumerate}[nosep]
  \item Compute the area of the ellipse $x^2/a^2 + y^2/b^2\leq1$ with
        the formula of \cref{stokes:ex:green}.
  \item Show that the integral of $\frac12(x\,\dd y - y\,\dd x)$ along
        the segment from $(x_1, y_1)$ to $(x_2, y_2)$ is
        $\frac12(x_1y_2 - x_2y_1)$. Deduce a formula for the area of a
        polygon from its vertices, admitting
        \cref{stokes:rem:corners}.
\end{enumerate}
\solutionlink{stokes:exr:3}
\end{exercise}

\begin{exercise}\label{stokes:exr:4}
On the torus obtained by rotating around the $z$-axis the circle of
radius $r$ centred at distance $R > r$ from the axis, with the
parametrization
$\psi(s, t) = \bigl((R + r\cos s)\cos t, (R + r\cos s)\sin t,
r\sin s\bigr)$:
\begin{enumerate}[nosep]
  \item compute the Gram matrix of the induced metric and the area of
        the torus;
  \item compute the integral of the function $h = z^2$ for the
        Riemannian volume.
\end{enumerate}
\solutionlink{stokes:exr:4}
\end{exercise}

\begin{exercise}\label{stokes:exr:5}
Let $R = [a, b]\times[c, d]$ and $\omega = P\,\dd x + Q\,\dd y$ smooth
on a neighbourhood of $R$. Prove directly, with Fubini's theorem, that
$\int_R\dd\omega$ is the sum of the integrals of $\omega$ along the
four sides, traversed counter-clockwise.
\solutionlink{stokes:exr:5}
\end{exercise}

\begin{exercise}\label{stokes:exr:6}
Let $D$ be a compact domain with smooth boundary in $\R^3$ and
$\sigma = i_E\,\mu$ as in \cref{stokes:ex:sphere}.
\begin{enumerate}[nosep]
  \item Show that $\lambda(D) = \frac13\int_{\partial D}\sigma$.
  \item Assume $0\notin\partial D$. Using
        \cref{differential-forms:exr:6}, show that
        $\int_{\partial D}\sigma/r^3$ is $4\pi$ if $0\in D$ and $0$
        otherwise.
\end{enumerate}
\solutionlink{stokes:exr:6}
\end{exercise}

\begin{exercise}\label{stokes:exr:7}
On the circle $S^1\subset\R^2$ let $\alpha$ be the restriction of
$x\,\dd y - y\,\dd x$.
\begin{enumerate}[nosep]
  \item Compute $\int_{S^1}\alpha$ for the two orientations of the
        circle, and $\int_{S^1}\abs\alpha$.
  \item Show that there is no function $f$ on the circle with
        $\dd f = \alpha$, and no $1$-form $\eta$ on the circle with
        $\int_{S^1}\eta\neq0$ that is exact.
\end{enumerate}
\solutionlink{stokes:exr:7}
\end{exercise}

\begin{exercise}\label{stokes:exr:8}
Let $p$ be the density of the Gaussian distribution $N(0, \Sigma)$ on
$\R^n$ and $X(x) = Ax$ for a matrix $A$.
\begin{enumerate}[nosep]
  \item Write Stein's identity for this $p$ and $X$, and check it with
        $\E[xx^\top] = \Sigma$.
  \item For the model $s(x) = -Bx$ with $B$ symmetric, compute
        $\E_p\bigl[\frac12\norm s^2 + \operatorname{div}s\bigr]$ and
        show that it is smallest for $B = \Sigma^{-1}$.
\end{enumerate}
\solutionlink{stokes:exr:8}
\end{exercise}

\begin{exercise}\label{stokes:exr:9}
For the graph of \cref{stokes:fig:graph}:
\begin{enumerate}[nosep]
  \item find a basis of $\ker\dual\dd$ and check the dimension given by
        \cref{stokes:prop:graph};
  \item decide whether $\alpha = (1, 2, 1, 3, 2)$ and
        $\alpha' = (1, 2, 2, 3, 2)$, listed in the order
        $e_1,\dots,e_5$, are of the form $\dd f$, and find $f$ when
        they are;
  \item compute $L$ and check that its rows sum to $0$.
\end{enumerate}
\solutionlink{stokes:exr:9}
\end{exercise}
\end{notechapter}
